% ==================================================
% Fundamental Theorem of Finite Calculus
% Discrete Calculus
% ==================================================

\subsection*{Fundamental Theorem of Finite Calculus}

\begin{theorembox}{Theorem (Fundamental Theorem of Finite Calculus)}
\begin{theorem}[Fundamental Theorem of Finite Calculus]
\label{thm:fundamental-theorem-finite-calculus}
Let \(F\) be a discrete antiderivative of \(f\), so that
\(\Delta F(n) = f(n)\) for all \(n\). Then
\[
\sum_{k=a}^{b-1} f(k) = F(b) - F(a).
\]
\hyperref[prf:fundamental-theorem-finite-calculus]{\textit{Go to proof.}}
\end{theorem}
\end{theorembox}

\begin{remark*}[Standard quantified statement]
\(\forall f,F : \mathbb{Z}\to\mathbb{R},\;
\forall a,b \in \mathbb{Z},\; a \le b:\;
(\forall n,\; \Delta F(n) = f(n))
\Rightarrow
\displaystyle\sum_{k=a}^{b-1} f(k) = F(b) - F(a)\).
\end{remark*}

\begin{remark*}[Predicate reading]
\[
\Delta F=f\Longrightarrow \sum_{k=a}^{b-1}f(k)=F(b)-F(a).
\]
\end{remark*}

\begin{remark*}[Interpretation]
This theorem is the discrete analogue of the Fundamental Theorem of
Calculus,
\[
\int_a^b f'(x)\,dx = f(b)-f(a).
\]
In continuous calculus, differentiation and integration are inverse
operations. In discrete calculus, the analogous inverse pair is the
difference operator \(\Delta\) and summation. Any sum can be evaluated
by finding a discrete antiderivative of the summand --- just as any
definite integral can be evaluated by finding an antiderivative.
The theorem reduces the problem of computing \(\sum_{k=a}^{b-1} f(k)\)
to the algebraic problem of finding \(F\) with \(\Delta F = f\). Combined
with the falling factorial rules developed in the next section, this
gives a systematic method for evaluating a wide family of sums in
closed form.
\end{remark*}

\begin{dependencies}
\begin{itemize}
  \item \hyperref[def:discrete-antiderivative]{Discrete Antiderivative}
  \item \hyperref[thm:telescoping-identity]{Telescoping Identity}
\end{itemize}
\end{dependencies}
